\section{Simulation Evaluation}
\label{sec:sim-eval}

Before evaluating on real protein data, we assess all models on
controlled simulations where the ground truth is known.
This establishes which models can detect which signals, and
identifies the data regimes (divergence, sequence length)
where differences are most informative.

\subsection{Simulation design}

\subsubsection{Trees from Pfam}

Simulations use phylogenetic trees derived from real Pfam families,
ensuring realistic topology and branch length distributions.
For each family in the held-out test set:
\begin{enumerate}
\item Parse the seed MSA and compute pairwise distances under LG08.
\item Build a neighbor-joining tree.
\item Optionally scale all branch lengths by a factor
  $s \in \{1.0, 2.0\}$ to explore longer-branch regimes.
\end{enumerate}

\subsubsection{Generating models}

Each fitted model (Section~\ref{sec:models-fitted}) serves as a
generating model, plus 2--3 randomly perturbed variants of each
(multiplicative noise on rates, Dirichlet noise on weights).
This tests sensitivity to parameter perturbation and avoids
overfitting the evaluation to a single parameter set.

Two families of generating models:

\paragraph{TKF92/LG08 (order-0 baseline).}
The standard TKF92 indel model with LG08 substitution.
No domain structure, no inter-column correlations.
Indel rates: $\insrate = 0.046$, $\delrate = 0.054$,
fragment extension $r = 0.68$ (typical Pfam values).
Root sequence length drawn from the geometric stationary distribution;
resampled until within 50\% of the observed MSA column count.

\paragraph{MixDom (order-1).}
The full Nested TKF model with trained parameters
(each fitted model from Section~\ref{sec:models-fitted},
plus 2--3 randomly perturbed variants to assess sensitivity).
Simulation uses \emph{direct generative sampling} from the exact
nested model, \textbf{not} from the distilled order-1 WFST.
Each simulated sequence carries fully labeled annotations:
domain boundaries and types, fragment boundaries and types,
and site class assignments.

\textbf{Root sequence:} Sampled from the stationary distribution of
the nested model.  The number of domains $\sim \text{Geometric}(\kappa_0)$
where $\kappa_0 = \insrate_0/\delrate_0$; domain types $\sim \domdist$.
Within each domain, fragment count $\sim \text{Geometric}(\kappa_d)$;
fragment types follow the Markov chain with initial distribution
$\fragdist_d$ and transition matrix $\ext^{(d)}_{fg}$.
Per-position site class $\sim \classdist_{d,f,c}$; residue $\sim \eqm^{(c)}$.
Root length is resampled until within 50\% of the observed MSA column count.

\textbf{Evolution along branches:} For each branch (parent $\to$ child)
at divergence time $t$:
\begin{enumerate}
\item \emph{Top-level TKF91} on domain blocks:
  determines which ancestor domains survive (M), are deleted (D),
  and where new domains are inserted (I), using $(\insrate_0, \delrate_0, t)$.
\item \emph{Domain type assignment:} Inserted domains $\sim \domdist$.
\item \emph{Per-domain TKF91} on fragment links within each domain:
  M/I/D link path using $(\insrate_d, \delrate_d, t)$ with fragment
  extension rate $=0$ (fragment extension is a column-level decision,
  not per-branch).
\item \emph{Fragment expansion:} New insertion links get fragment type
  from the Markov chain $(\fragdist_d, \ext^{(d)}_{fg})$.
  Matched links keep ancestor's fragment type.
\item \emph{Per-position substitution:} Each site evolves under its
  class's transition matrix $P(t) = \exp(Q^{(c)} t)$;
  inserted sites sample from $\eqm^{(c)}$.
\end{enumerate}
Progressive simulation from root to leaves produces the full true MSA
with known homology and full annotations at every position.

\textbf{Note on distilled models:}
The distilled order-1 WFST (from the Woodbury distillation,
Section~\ref{sec:distillation}) is used for \emph{inference}
(progressive alignment, beam search, tree composition), not for
simulation.  Using the distilled model for simulation would test
the distillation approximation rather than the generative model itself.

\subsubsection{Stratification}

Families are stratified by:
\begin{center}
\begin{tabular}{lll}
\textbf{Divergence} & \textbf{Mean pairwise identity} & \textbf{Expected signal} \\
\hline
Close & 50--70\% & Strong tree signal, weak order-1 contribution \\
Moderate & 30--50\% & Moderate tree signal \\
Divergent & 18--30\% & Weak tree signal, order-1 most valuable \\
\end{tabular}
\end{center}
crossed with alignment length (short: 50--100 columns, long: 100--200 columns)
and branch scale factor ($s = 1.0$, $s = 2.0$).

\subsection{Evaluation metrics}

For each (generating model $G$, reconstruction model $R$, family $F$),
we report:

\subsubsection{Order-1 signal strength}

The mutual information between adjacent positions in the distilled
singlet WFST, measured as:
\begin{equation}
\mathrm{KL}_{\text{order-1}} = \sum_{\anctok_{\text{prev}}} \pi(\anctok_{\text{prev}})
  \sum_{\anctok_{\text{cur}}} P(\anctok_{\text{cur}} | \anctok_{\text{prev}})
  \cdot D_{\mathrm{KL}}\!\left[
    P(\anctok_{\text{next}} | \anctok_{\text{prev}}, \anctok_{\text{cur}})
    \,\|\,
    P(\anctok_{\text{next}} | \anctok_{\text{cur}})
  \right]
\end{equation}
where $P(\anctok_{\text{next}} | \anctok_{\text{prev}}, \anctok_{\text{cur}})$
is the order-1 WFST transition probability and
$P(\anctok_{\text{next}} | \anctok_{\text{cur}})$ is its marginal
(averaged over $\anctok_{\text{prev}}$).
Units: nats per position.
An order-0 model has $\mathrm{KL}_{\text{order-1}} = 0$.

\subsubsection{Pairwise alignment log-likelihood}

For each cherry pair $(x, y)$ in the tree at divergence time $\tau$:
\begin{equation}
\ell_{\text{pair}} = \log P(x, y \mid R, \tau)
\end{equation}
computed via the forward algorithm on the distilled pair HMM from model $R$.
Normalized by alignment length.
This measures how well the pair HMM explains observed sequence pairs,
independent of tree structure.

\subsubsection{TreeVarAnc ELBO}

The evidence lower bound from TreeVarAnc (Section~\ref{sec:varanc})
on the true MSA:
\begin{equation}
\mathrm{ELBO}(n) = \text{TreeVarAnc ELBO after } n \text{ BP sweeps}
\end{equation}
We report $\mathrm{ELBO}(1)$ (approximately Felsenstein column LL
plus singlet prior, before BP iterates) and $\mathrm{ELBO}(\infty)$
(after BP convergence, typically 10 sweeps).

The \textbf{BP improvement} $\Delta_{\mathrm{BP}} = \mathrm{ELBO}(\infty) - \mathrm{ELBO}(1)$
measures the value of order-1 structure: how much does iterating
belief propagation with the WFST factors improve the bound?
$\Delta_{\mathrm{BP}} > 0$ indicates that the order-1 model captures
inter-column correlations not available to column-independent inference.

\subsubsection{Root reconstruction accuracy}

Given the true MSA (alignment known), predict the root residue at
each column using the TreeVarAnc posterior:
\begin{equation}
\hat{a}_c = \arg\max_a \, q_{(\text{root}, c)}(a)
\end{equation}
Accuracy = fraction of columns where $\hat{a}_c$ matches the true
root residue.

\subsubsection{Alignment accuracy}

The true MSA (from simulation) is ``unaligned'' to produce raw
leaf sequences.  These are then re-aligned using FSA
(finite-state aligner) with each model's pair HMM.
Alignment accuracy is measured by:
\begin{itemize}
\item \textbf{Sum-of-pairs (SP) score:}
  fraction of homologous residue pairs correctly identified.
\item \textbf{Total-column (TC) score:}
  fraction of columns where all residues are correctly aligned.
\end{itemize}

\subsubsection{Combined alignment + reconstruction}

For each reconstruction model $R$:
\begin{enumerate}
\item Unalign the simulated leaf sequences.
\item Re-align using FSA with model $R$'s pair HMM.
\item Run TreeVarAnc on the estimated MSA to reconstruct the root.
\item Score the reconstructed root against the true root:
  \begin{itemize}
  \item Matches: positions where predicted and true root agree.
  \item Mismatches: positions present in both but with different residues.
  \item Insertions: positions in predicted root not in true root.
  \item Deletions: positions in true root not in predicted root.
  \end{itemize}
\end{enumerate}
This end-to-end evaluation captures both alignment and reconstruction
quality, reflecting the full inference pipeline.

Also run TreeVarAnc on the \emph{true} MSA (alignment given)
to isolate reconstruction quality from alignment quality.

\subsubsection{Annotation accuracy}

Since the direct generative simulation produces fully labeled sequences,
we can evaluate the quality of inferred annotations:
\begin{itemize}
\item \textbf{Domain accuracy:} fraction of columns where the inferred
  domain type matches the true domain type.
\item \textbf{Fragment state accuracy:} fraction of columns where the
  inferred fragment state matches the true fragment state.
\item \textbf{Site class accuracy:} fraction of columns where the
  inferred site class matches the true site class.
\end{itemize}
These require the partition phylo-HMM posterior
(Section~\ref{sec:partition-recon}) which produces per-column
posteriors over domain, fragment state, and site class.

\subsubsection{Parameter recovery}

To validate the training pipeline, we simulate data from known
parameters, fit the model from random initialization, and measure
parameter recovery:
\begin{itemize}
\item \textbf{Rate recovery:} relative error in $(\insrate_d, \delrate_d)$.
\item \textbf{Domain weight recovery:} KL divergence between true and
  fitted $\domdist$.
\item \textbf{Extension matrix recovery:} Frobenius norm of
  $\ext^{(\dom)}_{\text{true}} - \ext^{(\dom)}_{\text{fitted}}$.
\item \textbf{Classdist recovery:} per-domain KL divergence between
  true and fitted $\classdist_{\dom,\srcfrag}$.
\end{itemize}

\subsection{Expected outcomes}

\begin{enumerate}
\item \textbf{Model matching:}
  When $G = R$ (reconstruction model matches generating model),
  the ELBO and pairwise LL should be highest.
  This validates that the evaluation metrics are sensitive to
  model fit.

\item \textbf{BP improvement under order-1 generation:}
  When $G$ is MixDom (order-1), $\Delta_{\mathrm{BP}}$ should be
  positive for MixDom reconstruction models and near-zero for
  Felsenstein.
  When $G$ is TKF92/LG08 (order-0), $\Delta_{\mathrm{BP}}$ should
  be near-zero for all reconstruction models.

\item \textbf{Divergence effect:}
  $\Delta_{\mathrm{BP}}$ should be largest for divergent families
  (weak tree signal), where inter-column information from the WFST
  has the most to contribute.

\item \textbf{Reconstruction accuracy:}
  MixDom models should outperform Felsenstein on MixDom-generated
  data, especially at high divergence.
  On TKF92-generated data, all models should perform similarly
  (no order-1 to exploit).
\end{enumerate}
